
\leaf{1}{
	$\lim_{x\to 0}\frac{\sin x\cdot\ln\braI{1+x}-x^{2}}{x\braI{1-\cos x}}=\underline{\qquad\qquad}$ .
}

\leaf{2}{
	$\Int{0}{1}{\frac{x\rme^{x}}{\braI{1+x}^{2}}}{x}=\underline{\qquad\qquad}$ .
}

\leaf{3}{
	设函数 $f\braI{x}=\braI{x^{3}+1}^{n}\arctan x$ , 则 $f\fnder{n}\braI{-1}=\underline{\qquad\qquad}$ .
}

\leaf{4}{
	设 $x=g\braI{y}$ 是可导函数 $y=f\braI{x}$ 的反函数, 并且 $f\braI{2}=3$ , $f\fder\braI{2}=\sfrac{1}{4}$ , 令 $F\braI{x}=f^{2}\braI{4x-g\braI{2x+1}\!}$ , 则 $F\fder\braI{1}=\underline{\qquad\qquad}$ .
}

\leaf{5}{
	由方程 $x^{3}+y^{3}=y^{2}$ 确定的函数 $y=f\braI{x}$ 的斜渐近线为 $\underline{\qquad\qquad}$ .
}

\bigskip

\stem{A}{
	\par 设函数 $f$ 在 $\bfR$ 上具有一阶连续导数, 并且 $f\braI{0}=f\fder\braI{0}=0$ , $f\fdder\braI{0}$ 存在, 记 $$g\braI{x}=\casesII{\frac{f\braI{x}}{x}}{x\ne 0}{-1}{a}{x=0}\ .$$
	\branch{A$_{\bm1}$}{
		已知 $g$ 为 $\bfR$ 上的连续函数, 求常数 $a$ 的值.
	}
	\branch{A$_{\bm2}$}{
		对于\textbf{A$_{\bm1}$}中确定的常数 $a$ , 讨论 $g$ 的可导性及其导函数的连续性.
	}
}

\stem{B}{
	\branch{B$_{\bm1}$}{
		已知函数 $f$ 在 $\braVII{0\and+\infty}$ 上连续且单调递增, 并且 $f\braI{x}>-1$ , 证明: $$F\braI{x}=\Int{0}{x}{\braI{2t-x}f\braI{x-t+\Int{0}{x-t}{f\braI{s}}{s}}}{t}$$ 在 $\braI{0\and+\infty}$ 上严格递减.
	}
}

\newpage

\stem{C}{
	\par 设函数 $f$ 是 $\braI{0\and 1}$ 上的连续函数.
	\branch{C$_{\bm1}$}{
		已知 $f\braI{x}\geqslant 0$ 并且 $$\Int{0}{1}{f\braI{x}}{x}=0\ ,$$ 求函数 $f$ 的表达式.
	}
	\branch{C$_{\bm2}$}{
		已知 $$3\Int{0}{1}{f^{2}\braI{x}}{x}+4=12\Int{0}{1}{xf\braI{x}}{x}\ ,$$ 求函数 $f$ 的表达式.
	}
}

\stem{D}{
	\par 设函数 $f$ 是 $\braII{0\and 1}$ 上的可导函数.
	\branch{D$_{\bm1}$}{
		已知 $f\braI{0}=f\braI{1}=0$ , 证明: 存在 $\xi\in\braI{0\and 1}$ 使得 $f\fder\braI{\xi}+f^{2}\braI{\xi}=0$ .
	}
	\branch{D$_{\bm2}$}{
		已知 $f\braI{0}=1$ 且 $f\braI{1}=\sfrac{1}{2}$ , 证明: 存在 $\xi\in\braI{0\and 1}$ 使得 $f\fder\braI{\xi}+f^{2}\braI{\xi}=0$ .
	}
}

\stem{E}{
	\branch{E$_{\bm1}$}{
		证明: 当且仅当 $a=1$ 时, 存在 $\braVII{0\and+\infty}$ 上的恒正连续函数 $f$ 使得对一切 $x>0$ 成立 $$\Int{0}{x}{f^{2}\braI{t}}{t}=\frac{1}{ax}\braI{\Int{0}{x}{f\braI{t}}{t}}^{2}\ .$$
	}
}

\newpage
